\documentclass[UTF8, a4paper, oneside]{ctexbook}

% --- 常用宏包 ---
\usepackage{amsmath, amssymb, amsfonts}
\usepackage{graphicx}
\usepackage{booktabs}
\usepackage{algorithm}
\usepackage{algpseudocode}
\usepackage{soul}           % 高亮文本
\usepackage{xcolor}         % 颜色支持
\definecolor{hlbg_y}{RGB}{255, 230, 153}  % 浅黄色高亮背景
\definecolor{hlbg_b}{RGB}{220, 235, 250}  % 浅蓝色
\definecolor{hlbg_g}{RGB}{220, 250, 230}  % 浅绿色

\usepackage{geometry}
\geometry{
    a4paper,
    left=3cm,
    right=3cm,
    top=2.5cm,
    bottom=2cm
}

% --- 图表标题格式设置 (核心修复区) ---
\usepackage{caption}
% 必须先声明，caption 才能识别 song 和 wuhao 这两个 key
\DeclareCaptionFont{song}{\songti}
\DeclareCaptionFont{wuhao}{\zihao{5}}

\captionsetup{
    font={song, wuhao}, 
    labelformat=simple,
    labelsep=quad,
    textformat=simple
}
% 统一应用，不再报 undefined 错误
\captionsetup[table]{font={song, wuhao}}
\captionsetup[figure]{font={song, wuhao}}

% --- 页眉页脚 ---
\usepackage{fancyhdr}
\pagestyle{fancy}
\fancyhf{}
\fancyhead[C]{\songti\zihao{-5} 中山大学本科毕业论文}
\fancyfoot[C]{\songti\zihao{-5} \thepage}

% --- 脚注格式：修正版 ---
\usepackage{footmisc}

% 1. 设置脚注正文的字体和字号（宋体小五）
\renewcommand{\footnotelayout}{\songti\zihao{-5}}

% 2. 修正脚注编号格式（非上标，且保持编号后的间距）
\makeatletter
\renewcommand\@makefnmark{\hbox{\normalfont\@thefnmark}} % 编号不出上标
\renewcommand\@makefntext[1]{%
    \noindent\makebox[1.5em][l]{\@thefnmark}#1% 编号左对齐，留出固定宽度
}
\makeatother

% --- 标题页信息 ---
\def\thesisTitle{基于语义标识符和大语言模型的POI序列推荐研究}
\def\thesisTitleEn{POI Sequence Recommendation Based on Semantic Identifier and Large Language Model}
\def\authorName{叶文浩}
\def\studentID{22330128}
\def\department{人工智能学院}
\def\major{人工智能专业}
\def\supervisor{刘威}
\def\dateYear{2026}
\def\dateMonth{4}

% --- 正文章节格式设置 ---
\ctexset{
    chapter={
        format=\centering\heiti\zihao{3},
        beforeskip=0pt,
        afterskip=20pt
    },
    section={
        format=\raggedright\heiti\zihao{4},
        beforeskip=20pt,
        afterskip=10pt
    },
    % 二级及以下节标题：宋体小四号加粗左对齐空两格
    subsection={
        format=\songti\zihao{-4}\textbf,
        indent=2\ccwd,
        beforeskip=15pt,
        afterskip=8pt
    },
    subsubsection={
        format=\songti\zihao{-4}\textbf,
        indent=2\ccwd,
        beforeskip=10pt,
        afterskip=6pt
    }
}

\usepackage{hyperref} % 建议放在大部分宏包之后

\begin{document}



\frontmatter
% 摘要 (请确保这些 .tex 文件存在)
% \include{chinese_abstract}
% \include{english_abstract}

% 目录
\tableofcontents
\listoffigures
\listoftables

% 正文
\mainmatter
\songti\zihao{-4}  % 全局正文：宋体小四号


% 第一章 绪论
\chapter{绪论}
\label{chap:intro}

\section{选题背景与意义}
\label{sec:background}

随着移动互联网的快速发展和智能终端的普及，基于位置的服务（Location-Based Service, LBS）已成为人们日常生活中不可或缺的一部分。\textbf{兴趣点推荐（Point-of-Interest, POI推荐）}是LBS的核心技术之一，旨在根据用户的历史行为和偏好，预测用户下一步可能访问的地点，并向用户推荐感兴趣的POI。POI推荐在大众点评推荐、地图导航、社交网络好友推荐等场景中具有广泛的应用价值。

与传统的商品推荐不同，POI推荐面临三个独特的挑战：

\begin{enumerate}
    \item \textbf{地理空间约束}：用户的移动行为受到地理空间的严格限制，偏好访问距离较近的POI，且地理位置呈现明显的距离衰减特性。
    \item \textbf{时空序列依赖}：用户的访问行为具有强烈的序列相关性，既受短期序列模式影响（如``午餐后去咖啡厅''），也受长期周期性影响（如``每周一上午去健身房''）。
    \item \textbf{数据稀疏性问题}：相比于用户访问的POI数量，未访问的POI数量呈指数级增长，导致正样本严重不足。
\end{enumerate}

传统的协同过滤方法通过挖掘用户-POI交互矩阵中的相似性进行推荐，但难以有效建模用户行为的时空序列特性。近年来，深度学习技术在POI推荐领域取得了显著进展。\textbf{循环神经网络（RNN）}能够捕捉用户签到序列中的时序依赖，但难以建模非连续位置之间的远程时空关联。\textbf{注意力机制}的出现使得模型能够直接建立序列中任意两个位置之间的关联，但现有方法仍然依赖于物品ID的离散嵌入表示，存在冷启动和泛化能力不足的问题。

\textbf{大语言模型（Large Language Model, LLM）}的崛起为推荐系统带来了新的机遇。LLM具备强大的语义理解和推理能力，能够理解用户行为的语义意图，将推荐任务建模为自然语言理解与生成任务。然而，将LLM直接应用于POI推荐面临严峻挑战：LLM的输入为自然语言token序列，而POI的表示是离散的整数ID，缺乏语义信息。如何将LLM的强大能力与POI推荐的独特问题相结合，成为当前研究的重要方向。

\textbf{语义标识符（Semantic Identifier, Semantic ID）}技术提供了一种潜在的解决方案。Semantic ID通过深度学习模型将POI的文本属性（如类别、名称、描述）和地理属性编码为连续的语义向量，再量化为离散的token序列。与随机分配的整数ID不同，Semantic ID具有语义相关性——语义相近的POI在ID空间中也相近。这使得Semantic ID能够：
\begin{itemize}
    \item 缓解冷启动问题：新POI只要有文本描述即可获得有意义的ID表示
    \item 支持序列建模：量化的token序列可直接用于自回归语言模型
    \item 融合多模态信息：可同时编码文本、地理等多模态属性
\end{itemize}

基于上述背景，本研究提出一套基于\textbf{残差量化变分自编码器（RQ-VAE）}的Semantic ID生成系统，并结合\textbf{QLoRA微调的大语言模型}进行POI序列推荐。通过RQ-VAE将POI的多维特征（类别、地理位置、时间、属性文本）编码为三层语义ID，并针对语义ID冲突问题设计了冲突解决策略。在此基础上，使用Trie约束解码确保生成语义ID的合法性和有效性。

\section{国内外研究现状和相关工作}
\label{sec:related}

\subsection{传统POI推荐方法}
\label{sec:traditional_poi}

早期的POI推荐方法主要基于协同过滤（Collaborative Filtering, CF）。协同过滤通过挖掘用户行为数据中的相似性进行推荐，包括基于用户的协同过滤和基于物品的协同过滤。然而，协同过滤方法难以有效建模用户行为的时空序列特性，在稀疏数据集上表现不佳。

\textbf{因子分解个性化马尔可夫链（Factorizing Personalized Markov Chains, FPMC）}\cite{rendle2010fpmc}由Rendle等人于2010年提出，首次将矩阵分解与马尔可夫链统一到同一框架中。FPMC的预测公式为：

\begin{equation}
\hat{r}_{u,i} = \mathbf{p}_u^\top \mathbf{q}_i + \sum_{j \in S_{u}^{t-1}} \mathbf{x}_j^\top \mathbf{y}_i
\label{eq:fpmc}
\end{equation}

其中第一项建模用户的静态偏好（矩阵分解部分），第二项建模用户行为的序列转移（马尔可夫链部分）。FPMC能够同时捕捉用户的长期偏好和短期序列模式，但其存在两点局限：（1）马尔可夫链假设下一个POI仅与前$k$个POI相关，难以建模长程依赖；（2）基于矩阵分解的隐向量表示缺乏语义信息，无法处理新用户或新物品的冷启动问题。

\subsection{基于深度学习的POI推荐方法}
\label{sec:deep_poi}

\textbf{时空循环神经网络（Spatial Temporal Recurrent Neural Networks, ST-RNN）}\cite{liu2016predicting}由Liu等人于2016年提出，在RNN的隐状态转换中引入时间和空间距离矩阵，能够在一定程度上捕捉用户行为的时空模式。ST-RNN仍面临两个关键局限：（1）基于RNN的序列建模难以有效捕捉非连续位置之间的远程时空关联；（2）简单的距离矩阵无法表达复杂的时空交互模式。

\textbf{DeepMove}\cite{feng2018predicting}（Feng et al., WWW 2018）针对长期周期性问题进行了改进。DeepMove设计了历史注意力机制，能够从用户过去数周乃至数月的长期历史轨迹中，通过注意力机制检索与当前上下文最相关的行为模式。这一工作标志着POI推荐研究开始关注长程依赖和周期性问题。

\textbf{时空注意力网络（Spatial-Temporal Attention Network, STAN）}\cite{luo2021stan}（Luo et al., WWW 2021）放弃了RNN架构，采用纯注意力机制（Self-Attention）直接建模轨迹中任意两点之间的时空关联。STAN的核心创新在于：（1）通过多模态注意力显式计算任意两点之间的时空距离（基于经纬度的哈弗辛距离和时间差）；（2）考虑了空间上的``三角不等式''约束。

STAN代表了基于深度学习的POI推荐在当时的最优效果，但其仍然依赖于离散ID的嵌入表示，存在语义信息缺失和冷启动问题。

\subsection{基于图神经网络的POI推荐方法}
\label{sec:gnn_poi}

\textbf{图神经网络（Graph Neural Network, GNN）}能够自然地建模用户-POI交互数据和POI之间的地理关系。该方向的代表性工作将用户-POI交互和POI地理邻近关系联合建模到异构图上，通过图神经网络进行信息传播和聚合。融合地理信息的图神经网络能够有效提升POI推荐的准确性，但仍存在以下局限：（1）依赖于交互数据中的拓扑结构，对新POI的地理信息利用不充分；（2）图神经网络的消息传递机制难以捕捉用户行为的高阶序列模式；（3）同样面临ID语义缺失问题。

\subsection{基于自监督学习的序列化推荐方法}
\label{sec:ssl_rec}

自监督学习通过在无标注数据上设计辅助任务来增强序列表示，已被广泛应用于序列化推荐。\textbf{S³-Rec}\cite{zhou2020s3rec}（Zhou et al., CIKM 2020）采用掩码语言模型（Masked Language Modeling, MLM）目标进行自监督预训练，通过多视角对比增强序列表示。S³-Rec在多个公开数据集上取得了当时最优效果，证明了自监督预训练对序列化推荐的有效性。然而，S³-Rec同样依赖于物品ID的离散嵌入，无法解决冷启动和语义泛化问题。

\subsection{基于大语言模型的推荐方法}
\label{sec:llm_rec}

近年来，\textbf{大语言模型（LLM）}在推荐系统领域展现出巨大潜力。LLM凭借其强大的语义理解、推理和生成能力，能够实现多轮对话式推荐、推荐解释生成、跨域知识迁移等任务。

\textbf{生成式检索推荐}\cite{rajput2023recommender}（Rajput et al., NeurIPS 2023）探索了使用生成式模型进行推荐的范式，通过将物品编码为离散token序列，实现了端到端的生成式检索。该工作首次提出Semantic ID技术，证明了将推荐任务建模为序列生成问题的可行性。

然而，将LLM直接应用于POI推荐面临以下挑战：（1）LLM的输入为自然语言token序列，而POI表示为离散整数ID，缺乏语义信息和连续性；（2）POI的地理坐标等数值特征难以直接用自然语言描述；（3）现有的POI推荐数据集规模难以满足LLM微调的需求。

\subsection{语义标识符技术}
\label{sec:semantic_id}

\textbf{Semantic ID}技术由Rajput等人在NeurIPS 2023的``Recommender Systems with Generative Retrieval''论文中首次提出。该技术通过深度学习模型将物品映射到语义向量空间，再量化为离散的token序列。与随机整数ID不同，Semantic ID具有语义相关性——语义相近的物品在ID空间中也相近。

该工作的核心贡献包括：（1）将物品编码为离散的token序列，实现端到端的生成式检索；（2）Semantic ID的语义相关性使得模型具备良好的泛化能力，能够处理新物品的冷启动问题；（3）为推荐系统与语言模型的结合提供了新的思路。

然而，现有的Semantic ID研究在以下方面存在不足：（1）主要关注商品推荐场景，对POI推荐的独特挑战（如地理空间约束）关注不足；（2）量化方法较为简单，未能充分利用多模态特征；（3）缺乏对语义ID分布质量和冲突率的系统分析。

\section{研究空白与问题定义}
\label{sec:gap}

综合上述分析，现有POI推荐方法存在以下研究空白：

\begin{enumerate}
    \item \textbf{LLM与POI推荐的结合问题}：如何将LLM的强大语义理解能力与POI推荐的独特挑战（地理空间约束、时空序列依赖）有效结合，仍是一个开放问题。
    \item \textbf{Semantic ID在POI场景的适用性}：现有Semantic ID研究主要针对商品推荐场景，对POI推荐的独特挑战（如地理空间约束）关注不足，特别是地理空间特征如何融入语义编码尚未充分探索。
    \item \textbf{语义ID分布质量评估}：现有研究缺乏对Semantic ID分布的系统评估指标体系，包括冲突率、覆盖率、分布均匀性等方面的定量分析。
    \item \textbf{语义ID约束解码问题}：当使用LLM生成Semantic ID时，如何确保生成的ID是合法有效的（对应真实存在的POI），需要有效的约束解码机制。
\end{enumerate}

针对上述研究空白，本文提出一套基于Semantic ID和LLM的POI序列推荐系统。

\section{本文主要研究内容}
\label{sec:content}

本文的主要研究内容包括以下几个方面：

\begin{enumerate}
    \item \textbf{基于RQ-VAE的语义ID生成框架}：采用残差量化变分自编码器将POI的多维特征编码为三层语义ID，融合类别嵌入、地理坐标、Fourier时间特征和属性文本嵌入。
    \item \textbf{基于顺序索引的冲突解决机制}：对冲突的语义ID附加顺序索引后缀进行消歧，与前三层码本保持一致的格式。
    \item \textbf{语义ID分布质量分析}：建立覆盖率、冲突率、信息熵、Gini系数等指标，为语义ID生成质量提供定量评估依据。
    \item \textbf{基于QLoRA微调的LLM序列推荐}：使用QLoRA技术对Qwen3-8B-Instruct大语言模型进行指令微调，实现POI序列推荐。
    \item \textbf{基于Trie的约束解码}：设计基于Trie数据结构的约束解码机制，确保LLM生成的语义ID合法有效。
\end{enumerate}

\section{本文的主要贡献}
\label{sec:contribution}

本文的主要贡献可以概括为以下几点：

\begin{enumerate}
    \item \textbf{提出基于RQ-VAE的语义ID生成框架}：融合类别、地理、时间、属性文本等多维特征，通过残差量化生成具有语义相关性的三层POI标识符。
    \item \textbf{设计基于顺序索引的冲突解决机制}：对冲突的语义ID附加顺序索引后缀进行消歧，提升语义ID的覆盖率和区分度。
    \item \textbf{建立语义ID分布质量评估指标体系}：提出覆盖率、冲突率、信息熵、Gini系数等指标，为语义ID生成质量提供定量评估依据。
    \item \textbf{构建基于QLoRA-LLM的POI序列推荐系统}：将语义ID与大语言模型结合，配合Trie约束解码，实现可解释的POI序列推荐。
    \item \textbf{在Yelp真实数据集上进行实验验证}：在Philadelphia城市子集上进行实验，验证所提方法的有效性。
\end{enumerate}

\section{论文结构}
\label{sec:structure}

本文共分为六章，各章内容安排如下：

\textbf{第一章 绪论}：介绍POI推荐的背景与意义，分析国内外研究现状，指出研究空白，阐述本文的主要研究内容和贡献。

\textbf{第二章 问题定义和代表性方法}：给出POI推荐和语义ID的形式化定义，介绍RQ-VAE残差量化、QLoRA微调等代表性方法的技术细节。

\textbf{第三章 基于RQ-VAE的语义ID生成与冲突解决}：详细阐述语义ID生成框架、冲突检测与解决算法，以及语义ID分布质量分析指标。

\textbf{第四章 基于QLoRA-LLM的POI序列推荐}：介绍指令微调数据的构建方法、QLoRA微调策略和Trie约束解码机制。

\textbf{第五章 实验与分析}：在Yelp数据集上进行实验，与基线方法进行对比，进行消融实验分析，验证所提方法的有效性。

\textbf{第六章 总结与展望}：总结本文的研究工作和主要贡献，并展望未来的研究方向。


% 第二章 问题定义和代表性方法
\chapter{问题定义和代表性方法}
\label{chap:problem}

\section{POI推荐问题定义}
\label{sec:poi_problem}

\subsection{用户-POI交互序列}
\label{sec:user_poi_sequence}

给定用户 $u$ 的历史签到序列：

\begin{equation}
S_u = [(poi_1, t_1), (poi_2, t_2), \dots, (poi_n, t_n)]
\label{eq:user_sequence}
\end{equation}

其中 $S_u$ 表示用户 $u$ 的历史交互序列，$poi_i$ 为第 $i$ 次交互的POI，$t_i$ 为对应的时间戳。

\subsection{Next-POI预测任务}
\label{sec:next_poi_task}

基于用户历史序列和上下文信息，预测用户下一个可能访问的POI：

\begin{equation}
poi_{n+1} = \arg\max_{poi \in \mathcal{P}} P(poi \mid S_u, context)
\label{eq:next_poi}
\end{equation}

其中 $\mathcal{P}$ 为所有POI的集合，$context$ 包含时空上下文信息。

\section{语义ID的定义与表示}
\label{sec:sid_definition}

本文提出将POI映射为离散的语义标识符（Semantic ID），使其在保留语义相关性的同时具备可枚举的离散表示。

\subsection{层次化量化结构}
\label{sec:hierarchical_quant}

语义ID采用多层量化编码结构，从粗到细逐层区分POI：

\begin{equation}
SID(poi) = \langle a_{p1} \rangle \langle b_{p2} \rangle \langle c_{p3} \rangle
\label{eq:sid_format}
\end{equation}

例如：$\langle a_{12} \rangle \langle b_{34} \rangle \langle c_{56} \rangle$

其中 $\langle a_X \rangle$、$\langle b_X \rangle$、$\langle c_X \rangle$ 分别对应P1、P2、P3三层量化索引。三层量化使得具有相似语义特征的POI在ID空间中彼此靠近，从而支持基于ID序列的推荐学习。

各层的设计意图如下：P1层捕获宏观特征（如地理大区、核心业态）；P2层细化到城市级别特征；P3层编码细粒度差异（如同一区域内不同类别）。三层码本相互独立，协同实现语义粒度的层次化表征。

\subsection{冲突处理原则}
\label{sec:collision_principle}

由于离散码本容量有限，多个POI可能被量化到相同的语义ID。此时需要附加额外信息进行消歧，确保每个POI拥有全局唯一的标识符。消歧策略将在第三章中详述。

\section{代表性方法概述}
\label{sec:related_methods}

本章概述与本文研究密切相关的几类代表性方法，包括序列推荐、语义编码与大语言模型推荐三个方向。

\subsection{序列推荐方法}
\label{sec:sequential_rec}

序列推荐旨在利用用户历史行为序列建模其动态偏好。早期代表性方法包括马尔可夫链模型（如FPMC~\cite{rendle2010factorifying}），通过对用户-物品交互矩阵进行分解并引入马尔可夫链建模序列转移。后续研究在此基础上引入注意力机制（如STAN~\cite{wang2021spatial}）以建模更复杂的序列依赖关系。

深度学习方法进一步推动了序列推荐的发展。RNN类方法（如GRU4Rec~\cite{hidasi2015session}）利用循环神经网络建模长序列依赖；Transformer类方法（如SASRec~\cite{kang2018self}）通过自注意力机制实现更高效的位置建模；时空注意力网络（如STAN）则在注意力基础上显式融合空间和时间信息。

上述方法均以预测下一个交互物品为直接目标，在隐式表示空间中运作，难以直接解释推荐理由。

\subsection{语义编码方法}
\label{sec:semantic_encoding}

语义编码方法的核心思想是将物品映射到低维语义向量或离散代码，以支持可解释的推荐推理。

VGAE~\cite{kipf2016variational}和VAE系列方法利用变分自编码器学习物品的连续语义表示；向量量化技术（如VQ-VAE~\cite{van2017neural}）进一步将连续表示离散化为紧凑码字，支持基于码本的物品检索。

近年来，研究者开始关注语义标识符（Semantic ID）的设计，旨在生成兼具语义相关性和离散可枚举性的物品代码。在推荐场景中，语义ID有望替代传统物品ID，支持基于ID序列的生成式推荐~\cite{zhou2024sid}。

本文借鉴上述思路，提出基于残差量化变分自编码器（RQ-VAE）的语义ID生成方法，并将其作为大语言模型推荐的输入序列。

\subsection{大语言模型推荐}
\label{sec:llm_rec}

大语言模型（LLM）凭借其强大的语义理解和推理能力，在推荐系统中展现了新的潜力。近期研究表明，LLM可在零样本或微调条件下完成推荐任务，包括next-item预测~\cite{heidari2024llm4rec}、序列推荐~\cite{liu2024llmseq}等。

指令微调（Instruction Tuning）是实现LLM推荐的主流范式：通过设计包含用户信息、历史行为和推荐目标的Prompt模板，将推荐任务转化为条件文本生成任务。代表性工作包括InstructRec~\cite{zhange2023instructrec}和LLM-Rec~\cite{liu2024llm4rec}等。

然而，将LLM直接应用于推荐仍面临挑战：传统物品ID缺乏语义信息，难以被LLM理解；直接使用物品属性文本作为标识则序列过长且检索效率低。为解决上述问题，本文提出基于语义ID的Prompt设计策略，并通过QLoRA微调使LLM适配POI推荐任务。

\section{数据划分策略}
\label{sec:data_split}

\subsection{Yelp数据集城市分布统计}
\label{sec:city_distribution}

Yelp数据集覆盖24个州、1,261个城市，共117,788个POI。城市POI数量分布呈长尾分布，如表所示：

\begin{table}[htbp]
\centering
\caption{Yelp数据集Top-10城市POI数量统计}
\label{tab:city_distribution}
\begin{tabular}{llr}
\toprule
排名 & 城市 & POI数量 \\
\midrule
1 & Philadelphia (PA) & 11,711 \\
2 & Tampa (FL) & 6,943 \\
3 & Tucson (AZ) & 6,701 \\
4 & Indianapolis (IN) & 6,010 \\
5 & Nashville (TN) & 5,554 \\
6 & New Orleans (LA) & 5,155 \\
7 & Edmonton (AB) & 4,381 \\
8 & Reno (NV) & 4,143 \\
9 & Saint Louis (MO) & 3,795 \\
10 & Santa Barbara (CA) & 2,826 \\
\bottomrule
\end{tabular}
\end{table}

\subsection{Philadelphia数据子集构建规则}
\label{sec:philly_data_rules}

当针对Philadelphia城市构建训练数据子集时，采用以下规则：

\textbf{1. 城市级数据抽取}：从按州/城市分层的目录结构（`PA/Philadelphia/`）中直接加载该城市的全部POI数据，避免跨城市数据污染。

\textbf{2. 特征维度设计}：采用GNPR-SID V2特征方案，总维度79维（无属性）或143维（含属性），具体构成为：
\begin{itemize}
    \item 类别嵌入：64维，使用SentenceTransformer或可学习哈希嵌入
    \item 空间坐标：3维，经纬度转换为单位球面3D坐标
    \item 时间特征：12维，Fourier展开捕捉周期性（6频率$\times$sin/cos）
    \item 属性嵌入（可选）：64维，从POI的attributes字段文本编码
\end{itemize}

\textbf{3. 数据质量过滤}：仅保留`is_poi=True`的记录，确保为有效商户而非其他类型实体。

\textbf{4. 时间分布聚合}：从用户签到记录（review_poi.json）中统计每个POI的访问时间分布，包括时段比例（早/午/晚/夜）和周内比例（工作日/周末）。

\textbf{5. 类别标签处理}：从POI的categories字段提取所有类别标签，映射到类别词表。单个POI最多保留5个类别，不足部分补零。

Philadelphia数据子集共11,711个POI，675,543条用户评论，为训练语义ID模型提供充足的样本量。

\subsection{选择Philadelphia作为训练数据的原因}
\label{sec:philly_selection}

\textbf{选择Philadelphia作为训练数据的原因}：
\begin{enumerate}
    \item \textbf{数据规模最大}：11,711个POI，675,543条用户评论（占全量数据集的14.9\%），数据量充足
    \item \textbf{避免跨城市复杂性}：单一城市区域，POI之间的地理距离和语义相关性更加一致
\end{enumerate}



% 第三章 基于RQ-VAE的语义ID生成与冲突解决
\chapter{基于RQ-VAE的语义ID生成与冲突解决}
\label{chap:rqvae_sid}

\section{多源特征编码器}
\label{sec:feature_encoder}

\subsection{特征构成}
\label{sec:feature构成}

POI的多维特征构成如表所示：

\begin{table}[htbp]
\centering
\caption{POI特征构成}
\label{tab:poi_features}
\begin{tabular}{lll}
\toprule
特征类型 & 维度 & 说明 \\
\midrule
category\_emb & 64 & 类别嵌入（SentenceTransformer或可学习） \\
spatial\_3d & 3 & 球坐标（lat/lon $\rightarrow$ 3D单位球面） \\
fourier\_time & 12 & Fourier时间特征（6频率 $\times$ sin/cos） \\
attr\_emb & 64 (可选) & Attributes文本嵌入 \\
\midrule
总计 & 79 (无attr) / 143 (有attr) & \\
\bottomrule
\end{tabular}
\end{table}

\subsection{类别编码器}
\label{sec:category_encoder}

类别编码器支持两种模式：
\begin{itemize}
    \item \textbf{预训练嵌入}：使用SentenceTransformer (all-MiniLM-L6-v2) 生成语义向量
    \item \textbf{可学习嵌入}：随机初始化，通过训练学习类别表示
\end{itemize}

多标签场景下，对所有类别嵌入取平均池化：

\begin{equation}
\mathbf{e}_{cat} = \frac{\sum_{j=1}^{N} \mathbf{e}_j \cdot \mathbb{1}_{c_j > 0}}{\sum_{j=1}^{N} \mathbb{1}_{c_j > 0}}
\label{eq:category_pooling}
\end{equation}

其中 $\mathbf{e}_j$ 为第 $j$ 个类别的嵌入向量，$c_j$ 为类别ID（大于0表示有效）。

\textbf{图示}：\figurename~\ref{fig:category_encoder}（待补充）

\subsection{空间坐标编码}
\label{sec:spatial_encoder}

将经纬度转换为3D球坐标，保持地理距离的几何特性：

\begin{equation}
\begin{aligned}
x &= \cos(\phi) \cdot \cos(\lambda) \\
y &= \cos(\phi) \cdot \sin(\lambda) \\
z &= \sin(\phi)
\end{aligned}
\label{eq:3d_sphere}
\end{equation}

其中 $\phi = \text{radians}(lat)$，$\lambda = \text{radians}(lon)$。

\textbf{图示}：\figurename~\ref{fig:spatial_encoding}（待补充）

该表示将地球表面映射到单位球面上，使得地理上相近的POI在向量空间中也相近。

\subsection{Fourier时间特征}
\label{sec:fourier_time}

编码访问时间的周期性模式（12维，6频率 $\times$ sin/cos）：

\begin{equation}
f_k^{sin} = \sin\left(\frac{2\pi k \cdot h}{24}\right), \quad
f_k^{cos} = \cos\left(\frac{2\pi k \cdot h}{24}\right), \quad k = 1, 2, \dots, 6
\label{eq:fourier_time}
\end{equation}

其中 $h$ 为访问时间的小时（0-24）。

Fourier特征能够有效捕捉时间序列的周期性规律，如``早高峰''、``午休''、``晚间活动''等模式。

\section{RQ-VAE语义ID生成框架}
\label{sec:rqvae_framework}

\subsection{模型架构}
\label{sec:rqvae_arch}

RQ-VAE的核心架构如图所示：

\begin{equation}
\text{Input} \xrightarrow{} \text{Encoder} \xrightarrow{} \text{RQ} \xrightarrow{} \text{Decoder} \xrightarrow{} \text{Output}
\end{equation}

RQ-VAE将多源特征编码器输出的特征向量作为输入，通过Encoder映射到64维的连续语义空间，然后通过残差量化器（RQ）将连续向量量化为三层离散的Semantic ID，最后通过Decoder重构原始特征向量进行验证。

\textbf{图示}：\figurename~\ref{fig:rqvae_arch}（待补充）

简化Encoder采用单层线性变换：

\begin{equation}
\mathbf{h} = \text{LayerNorm}\left(\text{GELU}\left(\mathbf{W}_e \mathbf{x} + \mathbf{b}_e\right)\right)
\label{eq:encoder_simple}
\end{equation}

其中 $\mathbf{x} \in \mathbb{R}^{79}$ 为输入特征，$\mathbf{h} \in \mathbb{R}^{64}$ 为编码后的隐向量。

原始方案采用多层MLP结构：
\begin{equation}
\text{Encoder}_{MLP} = [79, 512, 256, 128, 64]
\label{eq:encoder_mlp}
\end{equation}

RQ-VAE的主要超参数如表所示：

\begin{table}[H]
\centering
\caption{RQ-VAE核心参数配置}
\label{tab:rqvae_params}
\begin{tabular}{ll}
\toprule
参数 & 默认值 \\
\midrule
input\_dim & 79 \\
embedding\_dim & 64 \\
num\_quantizers & 3 \\
codebook\_size & 128 \\
commitment\_cost & 0.25 \\
\bottomrule
\end{tabular}
\end{table}

\subsection{残差量化器}
\label{sec:residual_quantizer}

残差量化器的核心机制如算法所示：

\begin{algorithm}[htbp]
\caption{残差量化算法}
\label{alg:rq}
\begin{algorithmic}[1]
\State $residual \gets z$ \Comment{输入向量 $z$（64维）}
\State $quantized\_sum \gets \mathbf{0}$ \Comment{初始化累加向量}
\For{$i \gets 1$ \textbf{to} $num\_quantizers$}
    \State $quantized\_i, indices\_i, loss\_i \gets quantizer\_i(residual)$ \Comment{第 $i$ 层量化}
    \State $quantized\_sum \gets quantized\_sum + quantized\_i$ \Comment{累加量化向量}
    \State $residual \gets residual - quantized\_i$ \Comment{计算残差传递到下一层}
\EndFor
\State \Return $quantized\_sum, [indices\_1, indices\_2, indices\_3]$ \Comment{返回量化向量和三组索引}
\end{algorithmic}
\end{algorithm}

\textbf{算法说明}：
\begin{itemize}
    \item 第1-2行：初始化残差向量和累加器
    \item 第3-7行：循环执行三层残差量化
    \item 第4行：每层使用独立码本将残差向量量化为128维码字之一
    \item 第5行：将当前层的量化结果累加到总量化向量
    \item 第6行：将残差减去量化值，传递给下一层继续量化
    \item 第8行：返回最终量化向量和P1/P2/P3三层索引
\end{itemize}

各层采用独立码本，码本大小为codebook\_size。

\textbf{EMA更新机制}：使用指数移动平均更新码本向量，死码（使用率低于平均值10\%）用随机样本重置。

\section{三层量化流程}
\label{sec:rq三层_quant}

RQ-VAE三层残差量化流程：

\begin{equation}
\text{输入向量 } z \text{ (64维)}
\end{equation}
\begin{equation}
\downarrow
\end{equation}
\begin{equation}
\text{P1量化}: z \rightarrow q_1, \quad r_1 = z - q_1
\end{equation}
\begin{equation}
\downarrow
\end{equation}
\begin{equation}
\text{P2量化}: r_1 \rightarrow q_2, \quad r_2 = r_1 - q_2
\end{equation}
\begin{equation}
\downarrow
\end{equation}
\begin{equation}
\text{P3量化}: r_2 \rightarrow q_3
\end{equation}
\begin{equation}
\downarrow
\end{equation}
\begin{equation}
\text{最终量化向量}: q_1 + q_2 + q_3
\end{equation}
\begin{equation}
\text{Semantic ID}: [idx_{P1}, idx_{P2}, idx_{P3}]
\end{equation}

其中P1层捕获宏观特征（地理大区、核心业态），P2层细化到城市级别特征，P3层编码细粒度差异。三层量化使得具有相似语义的POI在ID空间中相近。


\begin{figure}[H]
    \centering
    \includegraphics[width=0.75\linewidth]{figures/prefix1_semantic.png}
    \caption{第一层（P1）128个码字的Category与Region分布热力图}
    \label{fig:prefix1_semantic}
\end{figure}

\begin{figure}[H]
    \centering
    \includegraphics[width=0.75\linewidth]{figures/prefix2_semantic.png}
    \caption{Prefix=1条件下P2层86个码字的Category与Region分布热力图}
    \label{fig:prefix2_semantic}
\end{figure}

\begin{figure}[H]
    \centering
    \includegraphics[width=0.75\linewidth]{figures/prefix3_semantic.png}
    \caption{Prefix=1-11条件下P3层74个码字的Category与Region分布热力图}
    \label{fig:prefix3_semantic}
\end{figure}

\textbf{图3.1-3.3分析}：三层量化呈现明显的语义逐层细化规律。

\textbf{P1层宏观语义聚合}（图\ref{fig:prefix1_semantic}）：
\begin{itemize}
    \item \textbf{Category分布较为分散}：128个码字中，Top1纯度$\ge$90\%仅占11.7\%，72/128（56.3\%）的码字Top1占比低于70\%
    \item 高度混合码字多为跨业态聚合，如码字22（Health \& Medical 20.5\% + Arts \& Entertainment 18.2\%）、码字96（Fitness 22.6\% + Arts \& Entertainment 22.6\%），反映P1层捕获粗粒度地理-业态联合特征
    \item \textbf{Region分布高度碎片化}：90/128（70.3\%）的码字Top1邻里占比低于5\%，Region Top1均值仅4.7\%，说明P1层几乎不编码精确地理位置
\end{itemize}

\textbf{P2层城市级细化}（图\ref{fig:prefix2_semantic}）：
\begin{itemize}
    \item \textbf{Category纯度显著提升}：25/42（59.5\%）的码字达到Top1=100\%（完全纯净），高纯度（$\ge$90\%）占59.5\%
    \item 剩余混合码字多为相近品类共存，如码字1-34（Health \& Medical 50\% + Shopping 50\%）、码字1-3（Barbers 60\% + Beauty \& Spas 20\%）
    \item \textbf{Region分布集中度跃升}：11/42（26.2\%）的码字Top1邻里$\ge$50\%，Region Top1均值达46.5\%，表明P2层开始编码精确邻里级别的地理信息
\end{itemize}

\textbf{P3层细粒度消歧}（图\ref{fig:prefix3_semantic}）：
\begin{itemize}
    \item \textbf{Category纯度接近极值}：117/125（93.6\%）的码字Top1=100\%，仅8个码字存在品类共存，且均为近似品类对半分配（如Beauty \& Spas 50\% + Barbers 50\%、Health \& Medical 50\% + Counseling \& Mental Health 50\%）
    \item \textbf{Region分布高度集中}：93/125（74.4\%）的码字Top1邻里$\ge$50\%，Region Top1均值达86.2\%，表明P3层主要编码精确到街区的地理细分
    \item 少量P3码字仍存在Region Top1<50\%的情况，是因为同一街区内确实存在语义相近但实际不同的POI（如同一购物中心的不同店铺），此类冲突由冲突解决机制（$\langle d_{xx} \rangle$后缀）进一步消歧
\end{itemize}

\textbf{三层量化语义递增规律总结}：随着量化层级加深，Category语义纯度和Region地理集中度同步提升。P1层以粗粒度业态为主、地理为辅；P2层达到城市内邻里级别的语义-地理平衡；P3层聚焦到街区级精确语义，三层联合实现语义ID的层次化表征。

\section{损失函数与训练策略}
\label{sec:training_strategy}

\subsection{总损失函数}
\label{sec:rqvae_loss}

总损失函数定义为：

\begin{equation}
\mathcal{L}_{total} = \mathcal{L}_{recon} + \mathcal{L}_{quant} + \mathcal{L}_{align} + \mathcal{L}_{diversity}
\label{eq:rqvae_loss}
\end{equation}

其中：
\begin{itemize}
    \item $\mathcal{L}_{recon}$：重构损失，MSE(decoder\_output, input)
    \item $\mathcal{L}_{quant}$：量化损失，commitment\_loss
    \item $\mathcal{L}_{align}$：对齐损失，$1 - \text{cosine\_similarity}(quantized, encoded)$
    \item $\mathcal{L}_{diversity}$：码本均匀性损失（epoch $\geq$ 5后启用）
\end{itemize}

\textbf{设计说明}：
\begin{itemize}
    \item 重构损失确保量化后的向量仍保留原始特征信息
    \item 对齐损失确保量化向量与编码向量保持语义一致性
    \item 多样性损失延迟启用，避免训练初期干扰重构学习
\end{itemize}

\subsection{训练超参数配置}
\label{sec:training_hyperparams}

GNPR-SID V2方案下的完整训练超参数配置如下：

\begin{table}[htbp]
\centering
\caption{RQ-VAE训练超参数配置}
\label{tab:rqvae_training_params}
\begin{tabular}{ll}
\toprule
参数 & 默认值 \\
\midrule
embedding\_dim & 64 \\
num\_quantizers & 3 \\
codebook\_size & 128 \\
commitment\_cost & 0.25 \\
learning\_rate & 1e-3 \\
batch\_size & 32 \\
epochs & 100 \\
warmup\_epochs & 5 \\
diversity\_warmup\_epochs & 10 \\
diversity\_weight\_start & 0.05 \\
diversity\_weight\_end & 0.1 \\
use\_cosine & True \\
use\_ema & True \\
ema\_decay & 0.95 \\
\bottomrule
\end{tabular}
\end{table}

与原方案相比，GNPR-SID V2将输入特征维度从约256维降至79维，使编码器简化为单层线性变换（Linear(79$\rightarrow$64) + LayerNorm + GELU），大幅降低计算复杂度。

\subsection{训练收敛分析}
\label{sec:convergence_analysis}

在Philadelphia城市子集（9,256个POI，经k-core=5过滤）上训练100个epoch，训练曲线如下：

\begin{table}[H]
\centering
\caption{训练过程关键指标变化（PA\_main\_city\_cb128\_kcore实验）}
\label{tab:convergence}
\begin{tabular}{c|cc|c}
\toprule
Epoch & Train Loss & Val Loss & P3活跃码字 \\
\midrule
1   & 0.0890 & 0.0388 & 90/128 \\
5   & 0.0050 & 0.0041 & 110/128 \\
10  & 0.0020 & 0.0020 & 117/128 \\
50  & 0.0010 & 0.0010 & 128/128 \\
100 & 0.0009 & 0.0009 & 125/128 \\
\bottomrule
\end{tabular}
\end{table}

\textbf{收敛特征分析}：

\begin{itemize}
    \item \textbf{快速收敛}：模型在约5个epoch内即实现主要收敛（Train Loss从0.0890降至0.0050），约5个epoch即完成主要学习。
    \item \textbf{无过拟合}：验证损失（0.0009）与训练损失（0.0009）基本持平，100个epoch内未出现发散，说明模型泛化能力良好。
    \item \textbf{P3层码字动态调节}：P3层活跃码字在前50个epoch达到128/128（100\%），在epoch 100降至125/128（97.66\%），反映多样性损失逐渐将低频码字边缘化。
\end{itemize}

\subsection{动态多样性权重}
\label{sec:dynamic_diversity}

针对长尾分布数据的Warm-up策略：
\begin{itemize}
    \item 训练初期权重较小，允许模型先学习粗粒度聚类
    \item 逐渐增大权重，强制打散头部类别的拥堵
\end{itemize}

动态多样性权重的计算公式为：

\begin{equation}
w_{div}(e) =
\begin{cases}
w_{start} + \frac{e}{T_{warmup}} \cdot (w_{end} - w_{start}), & e < T_{warmup} \\
w_{end}, & e \geq T_{warmup}
\end{cases}
\label{eq:dynamic_diversity}
\end{equation}

其中 $e$ 为当前epoch，$T_{warmup}$ 为warmup轮数，$w_{start}=0.05$，$w_{end}=0.1$。

\textbf{设计说明}：PA\_main\_city实验中 diversity\_warmup\_epochs=10，即在前10个epoch内权重从0.05线性增长至0.1，此后保持0.1不变。实际效果显示P3层利用率在后期（epoch 100）从100\%降至99.22\%，说明多样性损失成功将3个低频码字逐渐淘汰。

\section{冲突解决机制}
\label{sec:collision_resolution_detail}

\subsection{冲突检测与解决}
\label{sec:collision_detection}

当多个POI被量化到相同的Semantic ID时，需要附加后缀进行消歧。冲突解决算法如下：

\begin{algorithm}[htbp]
    \caption{SID冲突解决算法}
    \label{alg:collision}
    \begin{algorithmic}[1] % [1] 表示显示行号
        \State $sid\_to\_pois \gets$ 按SID分组所有POI
        \For{$sid$ \textbf{in} $sid\_to\_pois$}
            \State $pois \gets sid\_to\_pois[sid]$
            \If{$|pois| = 1$}
                \State $resolved\_ids[pois[0]] \gets sid$
            \Else \Comment{多个POI共享相同SID}
                \For{$i \gets 0$ \textbf{to} $|pois|-1$}
                    \State $resolved\_ids[pois[i]] \gets sid + \langle d_{i+1} \rangle$
                \EndFor
            \EndIf
        \EndFor
        \State \Return $resolved\_ids$
    \end{algorithmic}
\end{algorithm}

\textbf{算法说明}：
\begin{itemize}
    \item 第1行：按Semantic ID对所有POI进行分组，得到 $\{SID \rightarrow [POI_1, POI_2, ...]\}$ 的映射结构
    \item 第2-7行：遍历每一个SID分组
    \item 第3行：获取当前SID分组下的所有POI列表
    \item 第4-5行：若该SID仅对应一个POI，则无需消歧，直接使用原SID作为最终标识符
    \item 第6-9行：若多个POI共享同一SID（发生冲突），则依次为每个POI附加顺序后缀 $\langle d_{i+1} \rangle$ 进行消歧
    \item 第10行：返回所有POI的最终解析标识符字典
\end{itemize}

冲突后缀格式为 $\langle d_{xx} \rangle$，与前三层码本保持一致的 $\langle X_{YY} \rangle$ 结构。

\textbf{设计说明}：早期曾考虑将PlusCode地理邻里编码注入Prompt以消歧，但可能引入噪声，后续简化为 $\langle d_{xx} \rangle$ 格式存储。

\subsection{冲突解决示例}
\label{sec:collision_example}

\begin{equation}
\text{原始SID}: \langle a\_{12} \rangle \langle b\_{34} \rangle \langle c\_{56} \rangle
\end{equation}
\begin{equation}
\text{冲突后}:
\begin{cases}
\langle a\_{12} \rangle \langle b\_{34} \rangle \langle c\_{56} \rangle \langle d\_1 \rangle & \text{第一个POI}\\
\langle a\_{12} \rangle \langle b\_{34} \rangle \langle c\_{56} \rangle \langle d\_2 \rangle & \text{第二个POI}
\end{cases}
\end{equation}

\section{语义ID分布质量分析}
\label{sec:sid_quality}

\subsection{核心指标体系}
\label{sec:sid_metrics}

语义ID分布质量评估指标如表所示：

\begin{table}[htbp]
\centering
\small
\caption{语义ID分布分析指标}
\label{tab:sid_metrics}
\begin{tabular}{lll}
\toprule
指标 & 说明 & 计算方法 \\
\midrule
覆盖率 (Prefix粒度) & 语义空间利用程度，唯一prefix数/总POI数 & $n_{prefixes} / n_{total}$ \\
无需后缀比例 & 一步到位比例，无需后缀的POI数/总POI数 & $n_{unique\_base} / n_{total}$ \\
Disambig率 & 需附加后缀的POI比例 & $n_{disambig} / n_{total}$ \\
碰撞率 (Collision Rate) & 有冲突的base SID占总SID比例 & $(n_{total\_tuples} - n_{unique\_tuples}) / n_{total\_tuples}$ \\
前缀分布熵 & 衡量prefix分布均匀程度，归一化为 $H/\log_2(n)$ & $H = -\sum p_i \log(p_i)$ \\
Gini系数（P1/P2/P3）& 衡量P1/P2/P3层codebook分布不均匀程度 & $G_{std} = \frac{\sum_{i}\sum_{j}|c_i-c_j|}{2n\sum_i c_i}$ \\
\bottomrule
\end{tabular}
\end{table}

\subsection{Philadelphia数据集实验结果}
\label{sec:philly_results}

在Philadelphia城市子集（11,711个POI）上训练100个epoch后的语义ID分布质量统计结果（来自PA\_main\_city实验组）如表所示：

\begin{table}[H]
\centering
\caption{Philadelphia数据集语义ID分布统计}
\label{tab:philly_sid_stats}
\begin{tabular}{llr}
\toprule
指标 & 数值 & 说明 \\
\midrule
总POI数 & 11,711 & - \\
覆盖率（prefix粒度） & 9,134 / 11,711 = 78.00\% & 语义空间利用程度\\
无需后缀比例 & 7,432 / 11,711 = 63.50\% & 一步到位，无需后缀的POI比例\\
Disambig率 & 36.5\% & 需附加后缀的POI比例\\
碰撞tuple数 & 1,702 & 有冲突的base SID数量\\
需消歧后缀POI数 & 4,279 & 冲突tuple中涉及的所有POI数量\\
前缀分布熵 & 12.99 / 0.99 & 原始熵 / 归一化熵 $H/\log_2(9134)$\\
Gini系数（P1/P2/P3）& 0.29 / 0.44 / 0.64 & P1/P2/P3层codebook分布不均匀程度\\
\midrule
\end{tabular}
\end{table}

\textbf{关键解读}：RQ-VAE三层量化中，P1/P2/P3层的标准Gini系数分别为0.29、0.44、0.64，显示层级越深分布越不均匀（残差头部效应加剧）。覆盖率78.00\%，无需后缀比例63.50\%，Disambig率36.5\%，前缀分布偏态分布是导致冲突的主要原因。相比早期方案的碰撞率，紧凑的特征（79维）大幅降低了语义空间坍缩风险。

\textbf{冲突规模分布}：2 POIs/tuple有1,187个，3 POIs/tuple有319个，4 POIs/tuple有114个，5 POIs/tuple有45个，6 POIs/tuple有18个，其余规模共28个tuple；最大冲突为16个POI共享同一SID。
\begin{figure}[H]
    \centering
    \includegraphics[width=0.75\linewidth]{figures/PA_main_city_collision_distribution.png}
    \caption{冲突规模分布（PA\_main\_city）}
    \label{fig:collision_distribution}
\end{figure}
\textbf{说明}：上图展示了冲突规模的分布情况，大多数冲突仅为2个ID。

\subsection{码本利用率分析}
\label{sec:codebook_usage}

在Philadelphia数据集上，三层量化码本的实际利用情况（来自PA\_main\_city实验组）：

\begin{table}[htbp]
\centering
\caption{码本利用率统计（PA\_main\_city实验，100 epochs）}
\label{tab:codebook_utilization}
\begin{tabular}{cccc}
\toprule
量化层 & 活跃码字数 & 码本大小 & 利用率 \\
\midrule
P1 & 125 & 128 & 96.88\% \\
P2 & 123 & 128 & 96.09\% \\
P3 & 105 & 128 & 82.03\% \\
\bottomrule
\end{tabular}
\end{table}

\textbf{分层利用率差异分析}：

\begin{itemize}
    \item \textbf{P1、P2层利用率高（96\%+）}：前两层活跃码字几乎占满128码本，表明粗粒度语义聚类充分，码本容量设计合理。
    \item \textbf{P3层利用率相对较低（82\%）}：第三层负责细粒度区分，当POI的语义特征差异已不足以在前两层被区分时，P3层面临更大的区分压力，导致部分码字使用率过低（接近死码阈值）而逐渐被EMA更新边缘化。训练结束时P3层死码数为3个（2.3\%），符合预期。k-core过滤后，P3层标准Gini从0.64降至0.43，残差分布偏态得到缓解。
    \item \textbf{死码重置机制生效}：代码中使用率低于平均值10\%的码字被标记为死码，用随机样本重置。该机制确保了码本的动态更新，避免低质量码字永久占据码本位置。
\end{itemize}

\subsection{多实验组对比分析}
\label{sec:multi_exp_comparison}

为验证GNPR-SID V2特征方案和消歧策略的有效性，本文在多个实验配置下进行对比：

\begin{table}[H]
\centering
\small
\caption{多实验组配置与效果对比}
\label{tab:multi_exp}
\resizebox{\textwidth}{!}{
\begin{tabular}{l@{\hspace{2pt}}|c|c|c|c}
\toprule
配置项 & PA\_main\_city & k-core(k=5) & PA全州 & attrs\_emb \\
\midrule
POI数量 & 11,711 & 9,256 & 26,938 & 11,711 \\
特征维度 & 79 & 79 & 79 & 143 \\
使用cosine loss & 是 & 是 & 是 & 是 \\
\midrule
P1利用率 & 96.88\% & 96.88\% & 100.00\% & 89.84\% \\
P2利用率 & 96.09\% & 95.31\% & 98.44\% & 82.03\% \\
P3利用率 & 82.03\% & 89.84\% & 92.97\% & 67.19\% \\
碰撞率 & 4.01\% & 1.83\% & 7.99\% & 4.69\% \\
Disambig率 & 36.5\% & 21.1\% & — & 41.1\% \\
覆盖率 & 78.00\% & 88.28\% & — & 74.34\% \\
\midrule
\end{tabular}
}
\end{table}

\textbf{对比结论}：

\begin{itemize}
    \item \textbf{数据质量过滤效果显著}：PA\_main\_city\_kcore在k-core=5过滤后（9,256 POIs），碰撞率从4.01\%降至1.83\%，Disambig率从36.5\%降至21.1\%，覆盖率从78.0\%升至88.3\%。说明被过滤的尾部长尾POI（评论稀疏的商户）是语义ID冲突的主要来源，过滤后模型聚焦于高质量样本，语义编码质量显著提升。
    \item \textbf{码本容量与冲突率呈负相关}：cb128（碰撞率1.83\%） vs cb64（碰撞率4.75\%），更大的码本容量能有效降低冲突率，验证了128码本对Philadelphia子集的适用性。
    \item \textbf{数据规模影响冲突率}：PA（26,938 POIs，7.99\%碰撞率） vs PA\_main\_city（11,711 POIs，4.01\%碰撞率），表明城市级子集因语义空间更集中，冲突率反而更低，验证了以Philadelphia作为训练子集的合理性。
    \item \textbf{attributes特征的作用有限}：PA\_philly\_attrs\_emb将特征维数扩展至143维（含384维attributes嵌入），但Disambig率（41.1\%）反而高于PA\_main\_city（36.5\%），P3层利用率骤降至67\%，说明高维稀疏的attributes文本嵌入在当前设置下未显著提升细粒度区分能力，反而因维度增加加剧了量化难度。注意该实验组同样使用了cosine loss，前述"cosine loss缺失导致利用率低"的结论不成立——低利用率的根本原因是特征维度过高。
\end{itemize}

\subsection{k-core 数据质量过滤消融实验}
\label{sec:kcore_ablation}

k-core过滤是提升语义ID质量的有效手段。为验证其作用，本节对比有无k-core过滤的两组实验结果（均使用128码本、cosine loss）：

\begin{table}[htbp]
\centering
\caption{k-core 数据质量过滤消融实验}
\label{tab:kcore_ablation}
\begin{tabular}{l@{\hspace{2pt}}|c|c}
\toprule
指标 & PA\_main\_city（无k-core） & k-core（k=5）\\
\midrule
训练POIs & 11,711 & 9,256 \\
Disambig率 & 36.5\% & 21.1\% \\
覆盖率 & 78.00\% & 88.28\% \\
碰撞率 & 4.01\% & 1.83\% \\
P1利用率 & 96.88\% & 96.88\% \\
P2利用率 & 96.09\% & 95.31\% \\
P3利用率 & 82.03\% & 89.84\% \\
Gini系数（P1/P2/P3） & 0.29 / 0.44 / 0.64 & 0.28 / 0.32 / 0.43 \\
冲突tuple数 & 47/1,172 & 17/927 \\
最大冲突规模 & 16 & 7 \\
\midrule
\end{tabular}
\end{table}

\textbf{对比结论}：

\begin{itemize}
    \item \textbf{碰撞率降低超过一半}：k-core过滤将碰撞率从4.01\%降至1.83\%，下降幅度达54\%，说明被过滤的尾部长尾POI（评论稀疏的商户）是语义ID冲突的主要来源。
    \item \textbf{Disambig率大幅改善}：从36.5\%降至21.1\%，意味着约63.5\%的POI可直接通过SID tuple映射，无需附加消歧后缀。
    \item \textbf{P3层利用率提升}：从82.03\%升至89.84\%，说明高质量样本的语义空间更易于被细粒度码本覆盖。同时P3层标准Gini从0.64降至0.43，残差分布偏态得到显著改善。
    \item \textbf{最大冲突规模显著缩小}：从16个POI共享同一SID降至7个，极端冲突情况大幅减少。
\end{itemize}

\subsection{各层码本热力图分析}
\label{sec:codebook_heatmap}

\begin{figure}[H]
    \centering
    \includegraphics[width=1.0\linewidth]{figures/PA_main_city_codebook_heatmap_combined.png}
    \caption{码本热力图（PA\_main\_city实验，三层对比）}
    \label{fig:codebook_heatmap}
\end{figure}

\begin{enumerate}
    \item \textbf{三层分布高度一致}：P1、P2、P3的热力图模式几乎相同，说明残差量化在各层稳定分摊量化任务。
    \item \textbf{码本利用充分}：几乎所有128个码字都在使用（颜色遍布全图），无明显空洞区域。
    \item \textbf{分布相对均匀}：颜色强度大多集中在中低范围，最大差异约6倍左右，无严重热点聚集，说明POI在语义空间中分布均衡，码本设计合理。
    \item \textbf{P3层热力略浅}：与P1、P2相比，P3层颜色整体偏浅（与82\%利用率相符），反映第三层部分码字承载量较低，且出现少数热点聚集，但可以接受。
\end{enumerate}

\section{本章小结}
\label{sec:chapter3_summary}

本章详细介绍了基于RQ-VAE的语义ID生成与冲突解决机制。主要内容包括：

\begin{enumerate}
    \item \textbf{多源特征编码器}：融合类别嵌入（64维）、3D球面坐标（3维）、Fourier时间特征（12维）等多种特征，总维度79维或143维。
    \item \textbf{RQ-VAE语义ID生成框架}：通过Encoder将多维特征映射到64维语义空间，然后通过三层残差量化生成离散的Semantic ID。
    \item \textbf{三层量化流程}：P1层捕获宏观特征，P2层细化到城市级别，P3层编码细粒度差异，各层独立码本。基于Philadelphia（11,711 POI）子集的实验表明，P3层Category纯度达93.6%，Region Top1均值达86.2%，三层联合实现语义ID的层次化表征。
    \item \textbf{损失函数与训练策略}：采用重构损失（+余弦对齐）、量化损失和动态调整的多样性损失联合优化。模型在约5个epoch内快速收敛，100个epoch内无过拟合迹象，CosSim稳定在0.994。
    \item \textbf{冲突解决机制}：采用网格邻里编码作为消歧后缀（\verb|[GG]|/\verb|[GG_n]|格式）。Philadelphia子集（11,711 POIs）全量统计Disambig率为36.5\%，经k-core=5过滤后（9,256 POIs）降至21.1\%，碰撞率从4.01\%降至1.83\%，验证了数据质量过滤对语义ID分布的显著改善作用。
    \item \textbf{语义ID分布质量分析}：多实验组对比验证了k-core过滤对降低冲突率的关键作用，以及attributes嵌入（384维）在当前设置下对区分能力提升有限。P3层利用率82\%，死码比例2.3\%，码本设计合理。
\end{enumerate}


% 第四章 基于QLoRA-LLM的POI序列推荐
\chapter{基于QLoRA-LLM的POI序列推荐}
\label{chap:qlora_llm}

\section{指令微调数据构建}
\label{sec:instruction_data}

\subsection{Prompt模板设计}
\label{sec:prompt_template}

本系统采用SID+时间历史的Prompt格式，将用户历史访问序列编码为自然语言提示。

\noindent\textbf{系统Prompt：}\fcolorbox{black}{hlbg_b}{%
\parbox{0.88\textwidth}{%
\small
"You are a POI next-visit prediction assistant. Predict only the next POI Semantic ID based on historical visit SIDs and timestamps. Output only one Semantic ID in angle-bracket format, such as <a\_12><b\_34><c\_56> or <a\_12><b\_34><c\_56><d\_1>."%
}}

\vspace{0.5ex}

\noindent\textbf{用户输入模板：}\fcolorbox{black}{hlbg_g}{%
\parbox{0.88\textwidth}{%
\small
"User\_\{user\_id\} checkin history: \{time1\} visited \{sid1\}, \{time2\} visited \{sid2\}, ... When \{target\_time\} user\_\{user\_id\} is likely to visit:"%
}}

该模板将用户历史访问序列（时间戳+Semantic ID）作为上下文，要求模型预测下一时刻可能访问的POI。

\subsection{训练数据生成流程}
\label{sec:data_generation}

训练数据构建流程如下：

\begin{algorithm}[htbp]
\caption{指令微调数据构建}
\label{alg:data_build}
\begin{algorithmic}[1]
\State 加载 Yelp 用户历史数据（reviews）和 POI 元数据
\State 加载预生成的 Semantic ID 映射
\For{\textbf{each} 用户 $u$}
    \State 按时间排序该用户的所有签到记录
    \State 构建访问序列 $visits = [(poi_1, t_1), (poi_2, t_2), \dots]$
    \For{$i \gets 1$ to $|visits|-1$}
        \State $history \gets visits[0:i]$  \Comment{取前$i$条作为历史}
        \State $target \gets visits[i]$  \Comment{第$i+1$条作为目标}
        \State 生成样本：$\{user\_id, history, target, target\_sid\}$
    \EndFor
\EndFor
\State 按 80\%/10\%/10\% 划分训练/验证/测试集
\end{algorithmic}
\end{algorithm}

数据格式化为对话格式：

\begin{equation}
\text{messages} = \left[
    \left(\text{"system"}, \mathbf{p}_{sys}\right),
    \left(\text{"user"}, \mathbf{p}_{user}(user\_id, history)\right),
    \left(\text{"assistant"}, sid_{target}\right)
\right]
\label{eq:chat_format}
\end{equation}

其中 $\mathbf{p}_{sys}$ 为系统提示，$\mathbf{p}_{user}$ 为用户输入模板，$sid_{target}$ 为目标Semantic ID。

\section{QLoRA微调策略}
\label{sec:qlora_config}

\subsection{量化配置}
\label{sec:quantization_config}

使用4bit NF4量化，配置参数如下：

\begin{itemize}
    \item \textbf{load\_in\_4bit}：启用4bit量化
    \item \textbf{bnb\_4bit\_quant\_type}：\verb|"nf4"|（NormalFloat4）
    \item \textbf{bnb\_4bit\_compute\_dtype}：\verb|bfloat16|
    \item \textbf{bnb\_4bit\_use\_double\_quant}：启用双重量化，进一步压缩
\end{itemize}

\textbf{图示}：\figurename~\ref{fig:qlora_quant}（待补充）

4bit量化显著减少显存占用，使得大模型微调可在消费级GPU上完成。

\subsection{LoRA配置}
\label{sec:lora_config}

LoRA微调的主要参数：

\begin{table}[htbp]
\centering
\caption{LoRA配置参数}
\label{tab:lora_config}
\begin{tabular}{ll}
\toprule
参数 & 默认值 \\
\midrule
r (秩) & 16 \\
lora\_alpha & 32 \\
target\_modules & q\_proj, k\_proj, v\_proj, o\_proj, \\
                 & gate\_proj, up\_proj, down\_proj \\
lora\_dropout & 0.05 \\
\bottomrule
\end{tabular}
\end{table}

LoRA通过在原模型权重旁添加低秩分解矩阵，仅训练这些新增参数，大幅降低微调成本。

\subsection{训练超参数}
\label{sec:training_params}

\begin{table}[htbp]
\centering
\caption{训练超参数}
\label{tab:training_params}
\begin{tabular}{ll}
\toprule
参数 & 默认值 \\
\midrule
max\_seq\_length & 704 \\
per\_device\_train\_batch\_size & 6 \\
gradient\_accumulation\_steps & 3 \\
learning\_rate & 2e-5 \\
num\_train\_epochs & 3 \\
warmup\_steps & 100 \\
bf16 & True \\
optim & paged\_adamw\_32bit \\
\bottomrule
\end{tabular}
\end{table}

有效batch size为 $6 \times 3 = 18$。

\section{Trie约束解码}
\label{sec:trie_decoding}

\subsection{TokenTrie数据结构}
\label{sec:trie_structure}

Trie树用于存储所有合法的Semantic ID序列，确保生成的ID对应真实存在的POI：

\textbf{图示}：\figurename~\ref{fig:trie_structure}（待补充）

\textbf{Trie节点结构}：
\begin{itemize}
    \item \textbf{children}：子节点映射，$\mathcal{C} = \{ (t_i \rightarrow node_i) \}$
    \item \textbf{is\_end}：是否为完整的Semantic ID终点
\end{itemize}

\textbf{核心操作}：
\begin{itemize}
    \item \textbf{add}$(\mathbf{t})$：将token序列 $\mathbf{t} = [t_1, t_2, \dots]$ 添加到Trie
    \item \textbf{traverse}$(\mathbf{t})$：沿token序列遍历Trie，返回当前节点
    \item \textbf{next\_tokens}$(\mathbf{t})$：返回当前前缀允许的下一个token集合 $\mathcal{N} \subset \mathcal{V}$
\end{itemize}

\subsection{TrieConstrainedLogitsProcessor}
\label{sec:trie_processor}

约束解码逻辑：

\begin{algorithm}[htbp]
\caption{Trie约束解码算法}
\label{alg:trie_decode}
\begin{algorithmic}[1]
\For{\textbf{each} $batch\_idx$} \Comment{遍历批次中的样本}
    \State $generated \gets input\_ids[prefix\_length:]$ \Comment{获取当前已生成的序列}
    \State $node \gets trie.traverse(generated)$
    \If{$node$ \textbf{is} None}
        \State $allowed \gets [EOS\_token\_id]$ \Comment{无效前缀，强制结束}
    \Else
        \State $allowed \gets node.children.keys()$ \Comment{获取Trie树下一层合法token}
        \If{$node.is\_end$ \textbf{and} $allow\_early\_stop$}
            \State $allowed.append(EOS\_token\_id)$
        \EndIf
    \EndIf
    \State $scores \gets mask\_disallowed(scores, allowed)$ \Comment{对非法token进行概率屏蔽}
\EndFor
\State \Return $scores$
\end{algorithmic}
\end{algorithm}

该机制确保：
\begin{enumerate}
    \item 生成的token序列始终是合法的Semantic ID前缀
    \item 若遍历失败（无效前缀），强制结束生成
    \item 支持early stop，当完整SID生成后允许EOS
\end{enumerate}

\section{实验设置}
\label{sec:experiment_setup}

\subsection{数据集}
\label{sec:dataset}

实验使用Yelp数据集的Philadelphia城市子集，该子集包含约14,569个POI和大量用户签到记录。

\begin{table}[htbp]
\centering
\caption{数据集统计}
\label{tab:dataset_stats}
\begin{tabular}{lr}
\toprule
项目 & 数值 \\
\midrule
POI数量 & 14,569 \\
用户数量 & ~10,000+ \\
签到记录数 & 100,000+ \\
\midrule
训练集比例 & 80\% \\
验证集比例 & 10\% \\
测试集比例 & 10\% \\
\bottomrule
\end{tabular}
\end{table}

\textbf{数据划分}：按时间顺序切分，确保测试集用户的历史完全包含在训练集中。

\subsection{评估指标}
\label{sec:eval_metrics}

采用以下评估指标：

\begin{table}[htbp]
\centering
\caption{推荐评估指标}
\label{tab:eval_metrics}
\begin{tabular}{lll}
\toprule
指标 & 说明 & 计算方法 \\
\midrule
HR@K & K次推荐中包含真实POI的比例 & $Hit@K / N_{total}$ \\
MRR@K & Mean Reciprocal Rank & $\frac{1}{K}\sum_{i=1}^{K}\frac{1}{rank_i}$ \\
NDCG@K & 归一化折现增益 & $DCG@K / IDCG@K$ \\
valid\_sid\_rate & 生成合法SID的比例 & $N_{valid} / N_{total}$ \\
\bottomrule
\end{tabular}
\end{table}

其中 $rank_i$ 为第 $i$ 个真实SID在Top-K预测中的位置。

\subsection{Baseline方法}
\label{sec:baseline_methods}

对比实验采用以下基线方法：

\begin{itemize}
    \item \textbf{FPMC}：因子分解个性化马尔可夫链，RecSys 2010
    \item \textbf{STAN}：时空注意力网络，WWW 2021
    \item \textbf{SARS}：序列感知推荐系统，RecSys 2020
\end{itemize}

\section{推荐效果对比实验}
\label{sec:recommendation_comparison}

\subsection{HR@K性能对比}
\label{sec:hrk_comparison}

\begin{table}[htbp]
\centering
\caption{HR@K性能对比}
\label{tab:hrk_comparison}
\begin{tabular}{lccc}
\toprule
方法 & HR@1 & HR@5 & HR@10 \\
\midrule
FPMC & - & - & - \\
STAN & - & - & - \\
SARS & - & - & - \\
\textbf{本文方法} & - & - & - \\
\bottomrule
\end{tabular}
\end{table}

\subsection{MRR@K性能对比}
\label{sec:mrrk_comparison}

\begin{table}[htbp]
\centering
\caption{MRR@K性能对比}
\label{tab:mrrk_comparison}
\begin{tabular}{lccc}
\toprule
方法 & MRR@1 & MRR@5 & MRR@10 \\
\midrule
FPMC & - & - & - \\
STAN & - & - & - \\
SARS & - & - & - \\
\textbf{本文方法} & - & - & - \\
\bottomrule
\end{tabular}
\end{table}

\subsection{NDCG@K性能对比}
\label{sec:ndcgk_comparison}

\begin{table}[htbp]
\centering
\caption{NDCG@K性能对比}
\label{tab:ndcgk_comparison}
\begin{tabular}{lccc}
\toprule
方法 & NDCG@1 & NDCG@5 & NDCG@10 \\
\midrule
FPMC & - & - & - \\
STAN & - & - & - \\
SARS & - & - & - \\
\textbf{本文方法} & - & - & - \\
\bottomrule
\end{tabular}
\end{table}

\section{消融实验}
\label{sec:ablation}

\subsection{冲突解决机制的作用}
\label{sec:collision_ablation}

\begin{table}[htbp]
\centering
\caption{冲突解决机制消融实验}
\label{tab:collision_ablation}
\begin{tabular}{lccc}
\toprule
配置 & HR@5 & MRR@5 & NDCG@5 \\
\midrule
无冲突解决 & - & - & - \\
有冲突解决 & - & - & - \\
\midrule
\textbf{提升} & - & - & - \\
\bottomrule
\end{tabular}
\end{table}

\textbf{分析}：冲突解决后缀 $\langle d_{xx} \rangle$ 使得具有相同基础SID的POI能够被区分，有效提升推荐准确性。

\subsection{Trie约束解码的作用}
\label{sec:trie_ablation}

\begin{table}[htbp]
\centering
\caption{Trie约束解码消融实验}
\label{tab:trie_ablation}
\begin{tabular}{lccc}
\toprule
配置 & Valid SID Rate & HR@5 & 生成延迟 \\
\midrule
无约束解码 & - & - & - \\
有约束解码 & - & - & - \\
\midrule
\textbf{提升} & - & - & - \\
\bottomrule
\end{tabular}
\end{table}

\textbf{分析}：Trie约束解码确保生成的Semantic ID对应真实存在的POI，有效避免无意义输出。

\section{生成式评测指标}
\label{sec:generation_metrics}

\begin{table}[htbp]
\centering
\caption{生成式评测指标}
\label{tab:gen_metrics}
\begin{tabular}{ll}
\toprule
指标 & 说明 \\
\midrule
exact\_match & Top-1与真实SID完全匹配的比例 \\
hit@K & 真实SID出现在Top-K的比例 \\
MRR@K & Mean Reciprocal Rank \\
valid\_sid\_rate & 生成的SID是合法POI的比例 \\
avg\_latency\_ms & 平均生成延迟（毫秒） \\
\bottomrule
\end{tabular}
\end{table}

\section{本章小结}
\label{sec:chapter4_summary}

本章介绍了基于QLoRA-LLM的POI序列推荐系统。主要内容包括：

\begin{enumerate}
    \item \textbf{指令微调数据构建}：设计SID+时间历史的Prompt模板，将用户历史访问序列编码为自然语言提示。
    \item \textbf{QLoRA微调策略}：采用4bit NF4量化和LoRA（r=16）微调Qwen3-8B-Instruct模型。
    \item \textbf{Trie约束解码}：通过Trie数据结构约束解码过程，确保生成的Semantic ID合法有效。
    \item \textbf{实验设置}：在Yelp Philadelphia数据集上进行实验，采用HR@K、MRR@K、NDCG@K等指标评估。
    \item \textbf{推荐效果对比}：与FPMC、STAN、SARS等基线方法进行对比。
    \item \textbf{消融实验}：验证冲突解决机制和Trie约束解码的有效性。
\end{enumerate}


% 参考文献
\clearpage
\begin{center}
    \heiti\zihao{3} 参考文献
\end{center}
\addcontentsline{toc}{chapter}{参考文献}
\bibliographystyle{gbt7714-numerical}
\bibliography{references}

% 致谢
\appendix
\include{acknowledgements}

% 附录
\include{appendix}

\end{document}